
%===================================
%===================================
\section{Conclusiones}
\label{sec:Conclusiones}

\begin{enumerate}

 \item Con 2 hilos y sin optimización, OpenMP solo alcanza un speedup promedio de 0.45x respecto a la versión secuelcial optimizada, es decir, es más del doble de lento. Con 4 hilos llega a 0.80x y recién a partir de 6 hilos (0.88x) empieza a acercarse a 1x. La causa es que el código paralelo de OpenMP sin optimización parte de un baseline serial lento, sin vectorización ni unrolling, y el paralelismo necesita varios hilos para compensar esa desventaja inicial.

    \item Con la optimización de compilador, OpenMP escala consistentemente: 1.63x a 4 hilos, 2.44x a 6, 3.25x a 8 y 4.94x a 12 hilos. El progreso es predecible, lo que indica que el overhead de OpenMP no interfiere de forma catastrófica con la vectorización del compilador. El máximo de 4.94x con 12 hilos es el mejor resultado absoluto entre todos los conteos de hilos de OpenMP.

    \item Los resultados muestran que los 6 hilos con Pthreads con optimización alcanza un speedup de 5.64x y la versión de 12 hilos de 5.47x, ambos por encima del mejor resultado de OpenMP (4.94x con 12 hilos). Una explicación posible es que las regiones paralelas de OpenMP introducen barreras implícitas y posibles problemas de aliasing de punteros que el compilador no puede resolver completamente, limitando la vectorización interna de cada hilo. En pthreads, el compilador tiene control total sobre cada función de hilo, lo que permite un uso más agresivo de AVX2.

    \item Con pthreads, tanto sin optimización como con optimización, 6 hilos superan a 12 hilos. Esto es coherente con la arquitectura del Ryzen 7600X, que cuenta con 6 núcleos físicos y 6 hilos SMT. En una carga fuertemente compute-bound como la multiplicación de matrices, el SMT compite por las unidades de punto flotante del mismo núcleo y no aporta beneficios. En cambio, con OpenMP, 12 hilos optimizados (4.94x) sí superan a 6 hilos optimizados (2.44x), lo que sugiere que el runtime de OpenMP maneja la afinidad o la distribución del trabajo de manera que logra extraer algún beneficio del SMT, aunque con menor eficiencia por hilo.

    \item La versión de 6 hilos de Pthreads con optimización, la eficiencia es $5.64 / 6 = 0.94$, prácticamente una escalabilidad lineal perfecta. En La versión de 12 hilos de Pthreads con optimización, la eficiencia cae a $5.47 / 12 = 0.46$, reflejando el impacto negativo del SMT. En la version de OpenMP con 12 hilos, la eficiencia es $4.94 / 12 = 0.41$, similar a pthreads con 12 hilos pero partiendo de un speedup menor. En la versión de OpenMP sin optimización, la eficiencia cae a $2.36 / 12 = 0.20$, lo que muestra que el overhead domina cuando el código serial es lento. El valor de 0.94 en la versión de Pthreads con 6 hilos optimizados es el dato más llamativo del benchmark y confirma que, a partir de 6 hilos, el cuello de botella principal es el hardware y no la implementación.

   \item Para $N=6400$ las tres matrices pesan aproximadamente 491MB, completamente fuera del alcance de cualquier caché en ambas máquinas. El problema viene del orden en que el algoritmo recorre los datos: en lugar de leer posiciones de memoria consecutivas para aprovechar la cacheline, salta de un extremo al otro de la matriz en cada paso. Eso hace imposible que el procesador anticipe qué dato va a necesitar a continuación, así que cada vez que lo requiere tiene que ir a buscarlo hasta la memoria principal, deteniéndose a esperar mientras el dato llega. Ese tiempo de espera es precisamente lo que los datos de perfilado muestran: la eficiencia del procesador, medida como instrucciones ejecutadas por ciclo, cae de 3.62 con $N=128$, donde todo el problema cabe cerca del procesador, a 2.63 con $N=1024$, cuando los datos ya no caben en los niveles más rápidos. Con $N=6400$ esa cifra sería aún menor, porque prácticamente cada acceso implica un viaje a la memoria principal.
   
   \item La optimización de compilador beneficia más a pthreads que a OpenMP.  
    Al pasar de sin optimización a con optimización, la implementación con pthreads de 6 hilos multiplica su speedup por 3.63x, pasando de 1.55x a 5.64x. En contraste, OpenMP con 12 hilos solo lo multiplica por 2.09x, de 2.36x a 4.94x. La diferencia es significativa y se explica porque el compilador tiene plena libertad para vectorizar las funciones de cada hilo en pthreads, al ser funciones C normales y aisladas. En cambio, dentro de una región paralela \texttt{\#pragma omp parallel for}, el compilador debe asumir aliasing de punteros entre hilos y la presencia de barreras implícitas de sincronización, lo que inhibe parte de la vectorización automática.

    \item Con optimización de compilador, pthreads con 6 hilos alcanza un speedup de 5.64x, mientras que OpenMP con 6 hilos solo llega a 2.44x usando los mismos recursos de hardware. Esto implica que pthreads es aproximadamente 2.3 veces más rápido. Esta diferencia no puede atribuirse únicamente al overhead de sincronización, que con 6 hilos es bajo, sino principalmente a las restricciones de vectorización presentes en OpenMP. Este caso actúa como el experimento de control más limpio del benchmark, ya que elimina los efectos del SMT y del uso de conteos de hilos desiguales.

    \item OpenMP con optimización escala casi linealmente al aumentar el número de hilos.  
    Dentro de la configuración de OpenMP con optimización, el speedup entre conteos consecutivos de hilos sigue una proporción cercana al ideal: de 2 a 4 hilos la mejora es 1.95x frente al ideal de 2x; de 4 a 6 hilos es 1.50x frente a 1.5x ideal; de 6 a 8 hilos es 1.33x frente a 1.33x ideal; y de 8 a 12 hilos es 1.52x frente a 1.5x ideal. Esta escalabilidad casi lineal se mantiene hasta 12 hilos, incluyendo el uso de SMT, y contrasta con pthreads, donde la ganancia colapsa al pasar de 6 a 12 hilos.

    \item El speedup aumenta con el tamaño del problema en todas las variantes paralelas.  
    En OpenMP con 12 hilos y optimización, el speedup crece desde 4.13x en $N=400$ hasta 5.80x en $N=6400$. Este patrón se repite en todas las configuraciones evaluadas. La razón es que el overhead de gestión de hilos es aproximadamente constante en tiempo absoluto, mientras que el trabajo útil crece como $O(N^3)$. A medida que el problema escala, la fracción paralela domina cada vez más el tiempo total de ejecución. Como consecuencia, los speedups promediados entre tamaños subestiman el rendimiento real en los problemas grandes, que son los más relevantes en la práctica.

    \item En pthreads sin optimización, pasar de 6 a 12 hilos produce una ligera caída de rendimiento, de 1.55x a 1.51x. En OpenMP sin optimización, en cambio, pasar de 8 a 12 hilos genera una ganancia real, de 1.67x a 2.36x. Esto sugiere que el runtime de OpenMP distribuye el trabajo y gestiona la afinidad de hilos de forma más eficiente que una implementación manual con pthreads cuando se entra en territorio SMT con código no vectorizado. No obstante, esta ventaja desaparece completamente cuando se habilitan las optimizaciones del compilador.

\end{enumerate}